\begin{abstract}
Weight-space learning --- predicting neural network properties directly from their parameters --- faces a fundamental challenge: permuting a network's neurons produces a functionally identical model with different raw weights. Non-invariant encoders treat these equivalent configurations as distinct, and the cost is measurable: a sinusoidal positional encoder (CISE) achieves $R^2 = -1.63$ when its predictions are averaged over functionally equivalent weight permutations, revealing that channel-position information is a primary predictive signal rather than background noise. We study this failure through a three-step causal framework --- architecture determines within-orbit representational variance (OrbitVar), OrbitVar propagates to permutation-induced prediction error ($\text{MSE}_\text{perm}$), and eliminating $\text{MSE}_\text{perm}$ improves downstream $R^2$ --- and close it empirically for the first time on ModelZooDataset CIFAR10-GS. Architecturally invariant encoders (DeepSets sum pooling) achieve OrbitVar at machine precision ($1.002 \times 10^{-14}$, twelve orders of magnitude below CISE), which reduces $\text{MSE}_\text{perm}$ from $3.35\times$ the total prediction error to essentially zero, yielding $R^2 = 0.9148$ --- a $+6.4$ percentage-point improvement over CISE with no additional training supervision. An unexpected entanglement between $\text{MSE}_\text{perm}$ and the residual error component reveals that non-invariant encoders and downstream predictors co-adapt, opening a new line of inquiry into the geometry of weight-space representations.
\end{abstract}
